% 09_references.tex —— 参考文献（统稿维护）
% 书目信息按 05_文献/文献映射.md 预填；标注 [check] 的字段定稿前从 PDF 首页核对
\begin{thebibliography}{29}

\bibitem{p1-03} Xia et al., Design of CO$_2$ huff-n-puff parameters for fractured tight oil reservoirs considering geomechanical effects. \emph{Processes}, 2024. %[check]
\bibitem{p1-04} Rezk and Foroozesh, Phase behavior and fluid interactions of a CO$_2$--light oil system at HPHT conditions. %[check]
\bibitem{p1-05} Li et al., The impact of nanoparticle adsorption on transport and wettability alteration. \emph{Frontiers in Physics}, 2019. %[check]
\bibitem{p1-06} Zhang et al., Quantitative investigation of nanofluid imbibition in tight oil reservoirs based on NMR. 2022. %[check]
\bibitem{p1-07} Wang and Sharma, Estimating unpropped fracture conductivity and fracture compliance from DFIT. %[check]
\bibitem{p1-08} Rostami et al., Shape factor for regular and irregular matrix blocks in fractured porous media. \emph{Petroleum Science}, 2019. %[check]
\bibitem{p1-09} Lee, Wheeler and Wick, Iterative coupling of flow, geomechanics and adaptive phase-field fracture. %[check]
\bibitem{p1-10} Hamiel, Lyakhovsky and Agnon, Poroelastic damage rheology. \emph{Geochemistry, Geophysics, Geosystems}, 2005. %[check]
\bibitem{p1-11} Jia, Ma, Cheng and Clarkson, Semi-analytical model for hydraulic fractures during flowback. \emph{Energy Science \& Engineering}, 2020. %[check]
\bibitem{p1-12} Hamiel, Lyakhovsky and Agnon, Coupled evolution of damage and porosity. \emph{Geophysical Journal International}, 2004. %[check]
\bibitem{p1-13} Zhang, Hu and Zhang, CO$_2$ huff-n-puff considering molecular diffusion and adsorption. \emph{Energies}, 2019. %[check]
\bibitem{p1-14} Merxhani, An introduction to linear poroelasticity. 2016. %[check]
\bibitem{p1-15} Nasser et al., Hybrid nanofluid wettability alteration and fracturing for CO$_2$ huff-n-puff. \emph{Energy Geoscience}, 2026. %[check]

\bibitem{p2-01} Rice and Cleary, Some basic stress diffusion solutions for fluid-saturated elastic porous media. \emph{Reviews of Geophysics}, 1976. %[check]
\bibitem{p2-02} Mahzari et al., Novel laboratory investigation of huff-n-puff gas injection for shale oils. %[check]
\bibitem{p2-03} Schmid and Geiger, Universal scaling of spontaneous imbibition. \emph{Water Resources Research}, 2012. %[check]
\bibitem{p2-04} Zhang et al., CO$_2$ huff-and-puff in tight sandstone: NMR study. %[check]
\bibitem{p2-05} Alotaibi et al., Stream and potential functions for transient flow. \emph{Fluids}, 2023. %[check]
\bibitem{p2-06} Fakcharoenphol et al., Effect of shut-in time on gas flow rate in hydraulically fractured shale reservoirs. \emph{Journal of Natural Gas Science and Engineering}, 2016. %[check]

\bibitem{p3-01} Wang et al., Optimization and analysis of CO$_2$ huff-n-puff using response surface methodology. \emph{Geofluids}, 2022. %[check]
\bibitem{p3-02} Zhang Jie et al., Sensitivity analysis and multiobjective optimization of CO$_2$ huff-n-puff. 2021. %[check]
\bibitem{p3-03} Kim and Durlofsky, A CNN--RNN proxy for robust optimization. %[check]
\bibitem{p3-04} Capolei et al., Risk minimization in life-cycle oil production optimization. %[check]
\bibitem{p3-05} Khanal and Shahriar, Optimization of CO$_2$ huff-n-puff in unconventional reservoirs. \emph{Energies}, 2023. %[check]

\bibitem{p4-01} Roberts et al., Cross-validation strategies for data with temporal, spatial, hierarchical, or phylogenetic structure. \emph{Ecography}, 2017. %[check]
\bibitem{p4-02} Saltelli et al., Variance based sensitivity analysis of model output: total sensitivity index. \emph{Computer Physics Communications}, 2010. %[check]
\bibitem{p4-04} Wang Qiang et al., Optimization method of post-fracturing soak time for shale oil reservoirs. \emph{Petroleum Exploration and Development}, 2022. %(in Chinese) %[check]
\bibitem{p4-05} Yang and Fang, Economic feasibility evaluation of CO$_2$ huff-and-puff. 2026. %[check]

\end{thebibliography}
